\section{Proposed Research Design}
\label{sec:design}

\subsection{Experimental Design}

The proposed study compares three map provenances:

\begin{enumerate}
\item a partisan-input-excluded geographic benchmark;
\item a commission-drawn map; and
\item an enacted legislative map.
\end{enumerate}

The information treatment has three levels: no process description, a
plain-language process description, and a description plus comparison. The
final registration must use a complete factorial design or state and justify
which cells are omitted. The earlier draft mixed a seven-cell layout with a
``$2\times3$'' label; that design is superseded.

The target sample is $N=2{,}400$ respondents. Final cell sizes will be fixed by
an a priori power analysis after the primary contrast, clustering unit,
multiple-testing family, and expected attrition are specified.

\subsection{Stimulus Governance}

Before registration, the study will publish:

\begin{itemize}
\item the legal provenance and effective date of every enacted or commission
map;
\item the complete benchmark profile, manifest, and assignment hash;
\item image-generation code, image hashes, dimensions, colors, labels, and
accessibility checks;
\item the state-selection rationale and any partisan-direction balancing rule;
and
\item cognitive-testing results for every process description.
\end{itemize}

The benchmark description must not say that the algorithm finds the ``most
compact'' map, guarantees fairness, or cannot be manipulated. It may state only
the inputs excluded from execution, the geographic objective, and the
reproducibility procedure.

\subsection{Recruitment And Ethics}

The proposed recruitment plan uses two independent panel sources with a target
of $1{,}200$ respondents each. Quotas, weighting, compensation, eligibility,
duplicate prevention, bot screening, and panel-specific exclusions must be
registered before launch.

No recruitment may begin until the project records the applicable
human-subjects ethics determination, informed-consent language, privacy plan,
retention period, and de-identification procedure.

\subsection{Registration}

The OSF or equivalent registration must freeze:

\begin{itemize}
\item hypotheses and confirmatory contrasts;
\item the primary and secondary outcomes;
\item sampling and stopping rules;
\item exclusions and missing-data treatment;
\item regression, clustering, weighting, and multiplicity procedures;
\item subgroup analyses;
\item stimuli and source hashes; and
\item the interpretation rules for null, heterogeneous, and adverse results.
\end{itemize}

No registration currently exists. A real URL and timestamp will replace this
statement only after the frozen package is public.

\subsection{Procedure And Exclusions}

The planned survey will collect demographics and baseline attitudes, present
the assigned stimulus under a minimum display time, collect map-specific
outcomes, administer attention and manipulation checks, and provide a
debriefing.

Exclusion thresholds will be fixed before launch. The final report will show
recruitment, exclusions, attrition, and analytic sample sizes by condition.
There are currently no recruited or excluded respondents.
